\section{What this doesn't do}
\label{sec:28.8}
It settles nothing. It is how the questions this book leaves open would be settled, and it will be slower than the labs for as long as they hold the weights, the compute and the people, which nothing in the design asks for and nothing in the budget should keep so: the chapter's precedents are CERN and the Office of Science's machines, not an arts endowment's grants, and a model a year behind the frontier is the least the floor needs, not the most the inquiry should want. It can't compel a lab outside its jurisdiction, or see into a classified system except through cleared teams whose findings come out by class. And it inherits the discourse's record: a discourse that turned objections into job titles and benchmarks can turn a public endowment into a grant category, and the measure that would show it is the same trend. What it changes is who holds the means. A finding that these systems can't be wronged would then be worth something, because a party that wanted the opposite answer could have looked with the same access; and a finding that they can be wouldn't depend on the permission of the party it costs.
